\chapter{Resultados y conclusiones}
Las arquitecturas comúnmente utilizadas en el procesamiento de imágenes requieren que los datos se encuentren ordenados en cierto formato, el cual puede ser en cuadrículas 3D, lo que permite que el proceso de distribución de peso se realice para optimizar el núcleo. Sin embargo, este enfoque puede resultar altamente costoso, ya que maneja una gran cantidad de parámetros que, 
en muchos casos, no aportan información relevante a la red y generan un desperdicio computacional significativo.
Para evitar estos inconvenientes, surge una arquitectura denominada PointNet, diseñada específicamente para procesar directamente nubes de puntos 3D sin la necesidad de 
transformarlas a una estructura regular como cuadrículas o vóxeles. PointNet aprovecha la naturaleza no estructurada de las nubes de puntos y utiliza operaciones simétricas, 
como el max pooling, para garantizar la invariancia al orden de los puntos. De este modo, permite realizar el procesamiento de datos 3D de manera más eficiente, 
capturando tanto características globales como locales directamente de los puntos de entrada, lo que simplifica y optimiza el análisis de formas y estructuras tridimensionales. \cite{pointnet}
En el presente capítulo analizaremos la forma en que el algortimo presentado en el capítulo anterior 
influye en los resultados de aplicar la arquitectura PointNet para la clasificación de nudos de Lissajous 
basado en la característica de Euler.
La descripción de la arquitectura está basada en el artículo \cite{pointnet}.

\section{Arquitectura PointNet}
La arquitectura, está diseñada de tal forma que recibe un 
conjunto de puntos sin orden en 3 dimensiones , es decir,
$\{P_i\ | i = 1, \cdots .,n\}$ donde cada punto $P_i$ es 
un vector con coordenadas $(x_i,y_i,z_i)$, es decir, una matriz $X \in \mathbb{R}^{n \times 3}$ la cual evita la pérdida
de información y reduce la complejidad computacional, a diferencia 
de otros algoritmos que convierten las nubes de puntos en imágenes 2D.

La entrada de la red se basa en tres propiedades principales de los puntos en el 
espacio Euclideano:

\begin{enumerate}
    \item \textbf{Sin orden}: La nube de puntos es un conjunto de puntos sin un orden
    en específico es decir la red consume  N puntos en 3d que son invariantes a $N!$ permutaciones en 
    el conjunto de datos de entrada.

    \item \textbf{Interacción entre puntos}: Los puntos provienen de un espacio métrico, lo cual
    otorga un subconjunto significativo de características por lo cual el modelo de poder capturar estructuras locales
    entre puntos cercanos.

    \item \textbf{Invariancia bajo transformaciones}: La nube de puntos es invariante bajo ciertas transformaciones
    por ejemplo, rotar y trasladar el conjunto de puntos no debería modificar la estructura global de dicho conjunto.
\end{enumerate}

\begin{figure}[h]
    \label{fig : PointNet Arquitecture}
    \includegraphics[width= \textwidth]{pointnet.png} 
    \centering
    \caption{Arquitectura PointNet \cite{pointnet}}
\end{figure}


En la medida en que el modelo debe permanecer invariante bajo permutaciones de los puntos de entrada, 
PointNet utiliza una función simétrica para agregar características de puntos individuales. 
Esta función simétrica es típicamente una operación de max pooling, que selecciona el valor máximo a lo largo de una dimensión específica.
La red se compone de varias capas de convolución 1D que operan sobre cada punto individualmente, seguidas de una capa de max pooling global que agrega las características extraídas de todos los puntos, pero principalmente sus capas siguen la siguiente ecuación que rige el funcionamiento y la validez de PointNet como aproximador universal de funciones continuas.
\\ Dado un conjunto de puntos sin orden $\{x_1,x_2, \cdots ,x_n\}$ con $x_i \in \mathbb{R}^d$ podemos definir
un conjunto de funciones que mapean el conjunto de vértices a un vector de la forma:

\[ \mathbf{x}_{i}^{\prime} = \mathbf{\gamma_{\Theta}} \left( \max_{j \in \mathcal{N}(i) \cup \{i\} } \mathbf{h_{\Theta}} (\mathbf{x}_j,p_j - p_i) \right) \]

Donde podemos recalcar, que cada punto se actualiza en función de la información proporcionada por sus vecinos $j \in \mathcal{N}(i)$.
De aquí que podemos notar que la función $\mathbf{h_{\Theta}} (\mathbf{x}_j,p_j - p_i)$ calcula la interacción entre el punto $\i$ y sus vecinos lo que puede considerarse como el mensaje que el punto $\i$ recibe.
Además la operción $\max$ indica que para cada punto $i$ se toma el valor máximo de las intereracciones de todos sus vecinos y de sí mismo actúando como un operador de agregación.
Por último notamos que la función $\gamma_{\theta}$ aplica una transformación no lineal a las representaciones agregadas, por lo tanto podríamos tomarlo como una función de actualización.
De esta forma podemos ver que la ecuación anterior la podriamos moldearcomo una GNN.


La arquitectura de PointNet se puede dividir en tres módulos principales:

\begin{enumerate}
    \item \textbf{Transformación de entrada}: PointNet aprende una transformación afín para alinear la nube de puntos de manera canónica, lo cual se logra 
    mediante uns subred llamada \textbf{T-Net} la cual predice una matriz de transofrmación $T \in \mathbf{R}^{3 \times 3}$, es decir \[ X^{*} = X \cdot T\], donde
    $X^{*}$ es la nube de puntos alineada y $T$ es la matriz predicha por \textbf{T-Net}
    \item \textbf{Extracción de características locales y globales}: A través de una red \text{MLP}, se extraen características locales de cada punto, es decir, para cada punto $x_i \in \mathbb{R}^3$,
     el MLP lo mapea a un espaciode caracteríticas de mayor dimensión: \[h_i = MLP(x_i) , \hspace{5pt} h_i \in \mathbb{R}^d\] 
     lo cual, al mapearse en todos los puntos, resulta en una matriz de características
    $H \in \mathbb{R}^{n \times d}$. Similar a la transformación de entrada, la red utiliza otra \textbf{T-Net} para aprender una transformación afín
    en el espacio de características, es decir: \[H* = H \cdot T_f\] donde $T_f \in \mathbb{R}^{d \times d}$ es la matriz de trasnformación aprendida.
    Luego, una capa de max pooling global agrega estas características en un vector de características globales $g \in \mathbb{R}^d$, es decir \[g_j = \max_{i = 1, \cdots ,n} H_{ij} \hspace{5pt} \forall j = 1, \cdots ,d \]
    \item \textbf{Clasificación o segmentación}: Dependiendo de la tarea, el vector de características globales se puede utilizar para la clasificación de la nube de puntos completa, o se pueden concatenar características locales y globales para realizar segmentación punto a punto.
    En el caso de la clasificación, podemos utilizar una red MLP para predecir la clase, es decir \[y = MLP(g)\] donde $y \in \mathbb{R}^k$ con $k$ clases y para el vector global $g$.

\end{enumerate}

\section{Resultados de la evaluación}

Ahora que describimos el funcionamiento y arquitectura de los modelos PointNet, presentaremos a continuación los resultados obtenidos de aplicar esta arquitectura combinado de el método de muestreo contra la arquitectura aplicada sin el método de muestreo. En esta sección, tabmién, se justifica el desempeño observado considerando las diferencias en la representación de datos y el procesamiento de la información.

\begin{figure}[h]
    \label{tab : Muestreo}
    \centering
    \begin{tabular}{| c | c | c | c | c |}
        \hline
         & precision & recall & f1-score & support \\ \hline
        0 & 0.85 & 0.85 & 0.85 & 20 \\ 
        1 & 0.77 & 0.50 & 0.61 & 20 \\ 
        2 & 0.95 & 0.90 & 0.92 & 20 \\ 
        3 & 0.83 & 0.95 & 0.88 & 20 \\ 
        4 & 0.80 & 0.80 & 0.80 & 20 \\ 
        5 & 1.00 & 0.30 & 0.46 & 20 \\ 
        6 & 0.57 & 1.00 & 0.73 & 20 \\ 
        7 & 0.72 & 0.90 & 0.80 & 20 \\ 
        8 & 0.84 & 0.80 & 0.82 & 20 \\ 
        9 & 0.90 & 0.90 & 0.90 & 20 \\ \hline
        accuracy & & & 0.79 & 200 \\
   macro avg & 0.82 & 0.79 & 0.78 & 200 \\ 
weighted avg & 0.82 & 0.79 & 0.78 & 200\\ \hline


    \end{tabular}
    \caption{Resultados obtenidos aplicando el muestreo general}
\end{figure}

\begin{figure}[h]
    \label{tab : sin_muestreo}
    \centering
    \begin{tabular}{| c | c | c | c | c |}
        \hline
         & precision & recall & f1-score & support \\ \hline
         0 & 0.32 & 1.00 & 0.49 & 20 \\
         1 & 0.29 & 0.35 & 0.32 & 20 \\
         2 & 0.58 & 0.95 & 0.72 & 20 \\
         3 & 1.00 & 0.30 & 0.46 & 20 \\
         4 & 0.00 & 0.00 & 0.00 & 20 \\
         5 & 0.42 & 0.55 & 0.48 & 20 \\
         6 & 0.00 & 0.00 & 0.00 & 20 \\
         7 &   1.00   &   0.25 &     0.40 &       20 \\
         8     &  0.61   &   0.85 &     0.71 &       20 \\
         9      & 0.81   &   0.65 &     0.72 &       20 \\ \hline

  accuracy & & &           0.49  &     200 \\
 macro avg   &    0.50  &    0.49  &    0.43   &    200 \\
weighted avg  &     0.50    &  0.49   &   0.43   &    200 \\ \hline


    \end{tabular}
    \caption{Resultados obtenidos sin aplicar el muestreo general}
\end{figure}

Podemos observar una diferencia notoria en lo resultados presentados, en primer lugar7 podemos observar que en el modelo con muestreo obtenemos métricas de precisión, recall y F1-score considerablemente más altos en la mayoría de las clases lo cual nos indica que este modelo clasifica correctamente las clases determinadas, también con base en el reccall podemos decir que el modelo detecta en la mayoría los casos positivos y por último la F1-score nos indica que la reducción de nodos ayudó a mantener la informacióin clave sin perder la representatividad en las clases.



\section{Conclusiones}